\subsection{Explicit sharp degrees and linear square tensor levels}

\begin{theorem}[DP58: explicit unit-parameter witnesses]\label{DP58}
Use the actual arbitrary-unit pair of DP46, with $p$ odd and $h\ge3$, and its primitive logarithmic splitting $\beta$ in additive coordinates and a normalized invariant integral frame.
For a primitive mixed line $(x,\mathbf y)=(ut,\mathbf vt)$, let $j$ be the first unit dual component and put
\[
 k=h-j,\qquad N=p^k,\qquad d=N+1,\qquad
 s=\frac1{p-1},\qquad
 w=\frac{sN}{N+1},\qquad c=\frac{s}{N+1}.
\]
For every unit $\lambda\in R^\times$ write
\[
 \exp\bigl(\lambda\beta(ut,\mathbf vt)\bigr)
 =\sum_n A_n(\lambda)t^n.
\]
At every degree $n_r=(p^k+1)p^r$, $r\ge0$, one has
\begin{align}
 v_p(A_{n_r}(\lambda))
 &=-\frac{p^{k+r}-1}{p-1},\label{eq:explicitkernelcoefficient}\\
 v_p(n_r!A_{n_r}(\lambda))
 &=\frac{p^r-1}{p-1}=cn_r-s.\label{eq:explicitkerneljet}
\end{align}
These valuations are independent of the unit $\lambda$.

For an eigenaxis $j$, put
$F_\lambda(x,y)=\exp(\lambda\beta(x,e_jy))$ and use the two-column polynomial $P_{a,b}$ of DP51.
Its value on $F_\lambda(x_1,y_1)F_\lambda(x_2,y_2)$ at zero satisfies
\begin{equation}\label{eq:explicitharmonicjet}
 v_p(H_{p^r,p^{k+r}}(\lambda))
 =\frac{p^r-1}{p-1}.
\end{equation}
For $j=1$ these are the source-shaped first-negative-row $\mathrm{HaO}$ polynomials; for other $j$ the assertion is only about the stated summed contraction unless a permitted row identification is supplied.

On the actual Jacobian of DP49 at $11$, every finite prefix of these explicit witnesses can be realized by global sections at square tensor powers satisfying
\[
 \kappa=M^2=An+O(\sqrt n),\qquad 11\nmid M,
\]
where $n$ is the largest selected degree and $A$ is a fixed interpolation constant from DP56.
The selected valuations include the source normalization $M^{-n}$, and the same tensor powers may be retained for number-field sections integral at the selected prime.
No section-height or prescribed-family conclusion is included here.
\end{theorem}

\begin{proof}
\emph{The full initial form of a critical exponential.}
Let $K'$ be a complete discretely valued extension with $v(p)=1$, containing $\mu$ with $v(\mu)=s$.
Suppose a zero-constant formal series $z$ has coefficients of valuation at least $s$ and its entire lowest-valuation form is
\[
 \overline{\mu^{-1}z}=u\mathbf Z^\alpha,
 \qquad u\ne0,\quad |\alpha|>0.
\]
Set $Q=\mu^{-1}z$.
The scalar in the $m$th term of
\[
 \mu^{-1}(\exp z-1)
 =\sum_{m\ge1}\frac{\mu^{m-1}}{m!}Q^m
\]
has valuation
\[
 (m-1)s-v_p(m!)=s(s_p(m)-1)\ge0.
\]
It is a unit precisely for $m=p^r$.
Reduction therefore gives the exact formal identity
\begin{equation}\label{eq:criticalexponentialinitial}
 \overline{\mu^{-1}(\exp z-1)}
 =\sum_{r\ge0}\epsilon_r u^{p^r}\mathbf Z^{p^r\alpha},
 \qquad
 \epsilon_r=\overline{\frac{\mu^{p^r-1}}{(p^r)!}}\ne0.
\end{equation}
Each coefficient before reduction is a finite sum because $Q$ has zero constant term.
The exponents on the right are distinct and their degrees tend to infinity.
Thus no convergence of an infinite sum in one fixed degree is being assumed, and no displayed leading coefficients cancel.
The hypothesis concerns the entire initial form; a single first ordinary-degree coefficient would not suffice.

\emph{Ordinary primitive lines.}
At the critical weight, the full initial-form calculation of DP46 gives a finite ramified scale $\varpi$ with $v(\varpi)=w$ and
\[
 S=\frac{N-1}{p-1},\qquad
 \mu=\frac{\varpi^{N+1}}{p^S},\qquad v(\mu)=s,
\]
such that
\[
 \overline{\mu^{-1}\beta(u\varpi T,\mathbf v\varpi T)}
 =\xi T^{N+1},\qquad \xi\ne0.
\]
This is the full initial form after the actual inverse logarithms; its scalar is the unit arrow product times
$(-1)^{k-1}\bar u\bar v_j^{p^k}$.
The invariant integral frame corrections are strictly higher by the invariant-frame part of DP46.
Apply \eqref{eq:criticalexponentialinitial} to $z=\lambda\beta(u\varpi T,\mathbf v\varpi T)$.
Every coefficient at degree $(N+1)p^r$ in the scaled exponential has valuation exactly $s$.
Hence
\[
 v_p(A_{n_r})=s-wn_r
 =-\frac{p^{k+r}-1}{p-1},
\]
which proves \eqref{eq:explicitkernelcoefficient}.
Since $n_r=p^{k+r}+p^r$ has digit sum two,
\[
 v_p(n_r!)=s(n_r-2),
\]
and addition gives \eqref{eq:explicitkerneljet}.
Every unit $\lambda$ multiplies the initial scalar by a nonzero residue, so the valuations are uniform over all such parameters.

\emph{The harmonic diagonal factor.}
With independent inputs, the same complete calculation gives
\[
 \overline{\mu^{-1}\lambda\beta(\varpi X,e_j\varpi Y)}
 =\xi_\lambda XY^N,\qquad \xi_\lambda\ne0.
\]
This is the two-variable initial-form calculation in DP51: the reduced scalar inverse is $X$, and the reduced last dual inverse is a unit times $Y^{p^{h-j-1}}$.
For $\widetilde F(X,Y)=F_\lambda(\varpi X,\varpi Y)$, the exponential identity yields
\[
 \overline{\mu^{-1}(\widetilde F-1)}
 =\sum_{r\ge0}\epsilon_r\xi_\lambda^{p^r}
       X^{p^r}Y^{Np^r}.
\]
Write $f_{a,b}=[x^ay^b]F_\lambda$ and set
\[
 D(Z)=\sum_{i\ge0}f_{i,i}Z^i,
 \qquad C(x,y)=F_\lambda(x,y)D(-xy).
\]
As $N>1$, the displayed initial form has no positive diagonal monomial.
Thus $v_p(\varpi^{2i}f_{i,i})>s$ for every $i\ge1$.
The values lie in one discrete value group, so the gap is uniform.
Consequently $\mu^{-1}(D(-\varpi^2XY)-1)$ is integral with zero reduction.
Multiplication by the integral series $\widetilde F$ preserves this zero reduction, and therefore
\[
 \overline{\mu^{-1}(C(\varpi X,\varpi Y)-1)}
 =\sum_{r\ge0}\epsilon_r\xi_\lambda^{p^r}
       X^{p^r}Y^{Np^r}.
\]
In particular $v_p(C_{p^r,Np^r})=s-w(N+1)p^r$.

The exact harmonic coefficient identity is
\begin{align*}
 H_{a,b}(\lambda)
 &=\sum_i(-1)^i\binom ai\binom bi
 (a-i)!(b-i)!(i!)^2 f_{a-i,b-i}f_{i,i}\\
 &=a!b!\,[x^ay^b]C(x,y),
\end{align*}
because the factorial factor in every summand equals $a!b!$.
This includes all the alternating terms.
At $(a,b)=(p^r,Np^r)$ both digit sums are one, so
\[
 v_p(a!b!)=s(a+b-2).
\]
Together with the coefficient of $C$ just computed, this proves \eqref{eq:explicitharmonicjet}.
The row-degree ratio is exactly $1:N$.

\emph{Linear square tensor powers.}
On the actual Jacobian at $11$, let $A\ge1,B_0\ge0$ be fixed constants such that every $n$-jet extends in every tensor power at least $An+B_0$, by DP56.
Put $m_0=\lceil\sqrt{An+B_0}\rceil$, choose $M=m_0$ unless $11\mid m_0$, and in that case choose $M=m_0+1$.
Then
\[
 11\nmid M,
 \qquad An+B_0\le M^2
 <An+B_0+4\sqrt{An+B_0}+4.
\]
At the chosen power $\kappa=M^2$, interpolate the algebraic prefactor to $U^{-\kappa}$ through degree $n$, as in DP54.
The normalized germ then has the pure kernel jet with unit exponent $\lambda=-\kappa$.
The exact formulas above hold for every unit exponent, so no further congruence on $M$ is necessary.
The source factor $M^{-n}$ is a unit.

The explicit ordinary degrees at $11$ are $12\cdot11^r$ and $122\cdot11^r$, each with derivative valuation $(11^r-1)/10$.
For the first-row harmonic family they are $(a,b)=(11^r,121\cdot11^r)$, with that same valuation.
Interpolating to the largest degree of any finite prefix realizes every earlier witness as well, with the same section and unit exponent.
For two columns use two copies of that section.
DP54's number-field approximation preserves all these finitely many nonzero values in valuation at the same chosen tensor power.
\end{proof}

This theorem concerns pure primitive kernels and global sections with deliberately prescribed finite prefactor jets. It does not assert these exact-degree values for every integral prefactor, does not realize infinitely many witnesses with one fixed section, and does not control the sections' archimedean norms or identify the prescribed translated family. The primitive projections and unit tensor parameter remain essential hypotheses; the nonunit regime of DP50 is separate.
